\chapter{Localization}
\label{ch:localization}
What localization is
\red{diagram of whole architecture}
\section{Testing Environment}
\section{Monte Carlo Localization}
\label{sec:mcl}

Monte Carlo localization (MCL) is a particle filter approach to localization that has
become one of the most popular localization algorithms in robotics~\cite{thrun2005probabilistic}.
It is easy to implement and works well across a wide range of localization problems. The
goal of localization is to estimate the rover's pose $x_t = (x, y, \theta)^\top$ in a known
map $m$, given every control $u_{1:t}$ and measurement $z_{1:t}$ received up to time $t$.
Equation~\ref{eq:belief} defines this estimate, called the belief, as a probability
distribution over all possible poses.
\begin{equation}
  \label{eq:belief}
  \mathrm{bel}(x_t) = p(x_t \mid z_{1:t}, u_{1:t}, m)
\end{equation}

MCL does not store the belief as a closed-form distribution. Instead, it keeps a set of $M$
samples, called particles, where each particle is one guess of the rover's pose.
Equation~\ref{eq:particle-set} shows the particle set at time $t$.
\begin{equation}
  \label{eq:particle-set}
  \mathcal{X}_t = \left\{ x_t^{[1]}, x_t^{[2]}, \dots, x_t^{[M]} \right\}
\end{equation}
Areas of the map that hold many particles are areas where the rover is likely to be, and
areas that hold few particles are areas where it is unlikely to be.

MCL was chosen over Gaussian filters such as the extended Kalman filter (EKF) and the
unscented Kalman filter (UKF) for several reasons~\cite{thrun2005probabilistic}. First, the
EKF and UKF represent the belief as a single Gaussian, so they can only track one pose
hypothesis at a time. They cannot solve global localization, where the starting pose is
unknown, and they cannot recover once the estimate has diverged. A particle set can hold
several hypotheses at once in different parts of the map. Second, the EKF linearizes the
motion and measurement models. The UKF avoids linearization but still assumes Gaussian
noise. MCL places no such restriction on either model. Third, EKF localization usually works
on a small set of extracted landmarks and must decide which landmark each measurement
belongs to. MCL can score raw range measurements directly against the map, which suits the
Cyclops rover, whose measurement is a dense range scan obtained from monocular depth
(Section~\ref{sec:mde-to-range}). Finally, the cost of MCL grows linearly with the number of
particles $M$, so it can be tuned to fit the computational budget of the rover.

Algorithm~\ref{alg:mcl} lists the MCL procedure, adapted from~\cite{thrun2005probabilistic}.
The particle index is written as $k$ instead of $m$ so that it is not confused with the
map $m$.

\begin{algorithm}[htbp]
  \caption{Monte Carlo Localization, adapted from~\cite{thrun2005probabilistic}}
  \label{alg:mcl}
  \begin{algorithmic}[1]
    \Function{MCL}{$\mathcal{X}_{t-1}, u_t, z_t, m$}
      \State $\bar{\mathcal{X}}_t \gets \emptyset$, \quad $\mathcal{X}_t \gets \emptyset$
      \For{$k = 1$ \textbf{to} $M$}
        \State $x_t^{[k]} \gets \Call{sample\_motion\_model}{u_t, x_{t-1}^{[k]}}$
        \State $w_t^{[k]} \gets \Call{measurement\_model}{z_t, x_t^{[k]}, m}$
        \State $\bar{\mathcal{X}}_t \gets \bar{\mathcal{X}}_t \cup \left\{ \langle x_t^{[k]}, w_t^{[k]} \rangle \right\}$
      \EndFor
      \For{$k = 1$ \textbf{to} $M$}
        \State draw $i$ with probability $\propto w_t^{[i]}$
        \State add $x_t^{[i]}$ to $\mathcal{X}_t$
      \EndFor
      \State \Return $\mathcal{X}_t$
    \EndFunction
  \end{algorithmic}
\end{algorithm}

Each call to Algorithm~\ref{alg:mcl} updates the belief from time $t-1$ to time $t$. The
inputs are the previous particle set $\mathcal{X}_{t-1}$, the latest control $u_t$, the
latest measurement $z_t$, and the map $m$. Line~2 creates two empty sets: a temporary set
$\bar{\mathcal{X}}_t$ that holds weighted particles, and the output set $\mathcal{X}_t$. The
rest of the algorithm runs in three steps.

\textbf{Prediction (line 4).} Each particle is moved forward by drawing a new pose from the
motion model, given the control $u_t$ and the particle's previous pose. Because the draw is
random, particles that started at the same pose spread apart, which reflects the uncertainty
in the rover's motion. The motion model used by the Cyclops rover is described in
Section~\ref{sec:motion-model}.

\textbf{Weighting (lines 5--6).} Each moved particle is given an importance weight
$w_t^{[k]}$ by the measurement model. The weight is high when the measurement $z_t$ the rover
actually received matches what it would expect to see from pose $x_t^{[k]}$ in the map $m$,
and low otherwise. Line~6 stores the particle together with its weight in
$\bar{\mathcal{X}}_t$. The measurement model used by the Cyclops rover is described in
Section~\ref{sec:measurement-model}.

\textbf{Resampling (lines 8--10).} A new set of $M$ particles is drawn with replacement from
$\bar{\mathcal{X}}_t$. Equation~\ref{eq:resample-prob} gives the probability that particle
$i$ is drawn on each draw, which is its weight normalized over all particles.
\begin{equation}
  \label{eq:resample-prob}
  P(\text{draw } i) = \frac{w_t^{[i]}}{\sum_{j=1}^{M} w_t^{[j]}}
\end{equation}
Particles with high weight are likely to be drawn several times, while particles with low
weight are likely to disappear. After resampling, the particles gather in the parts of the
map that agree with the measurement.

Line~12 returns $\mathcal{X}_t$, which represents the belief $\mathrm{bel}(x_t)$ and becomes
the input to the next call. A single pose estimate for the rover is taken as
the mean of the particle set \red{[confirm: mean vs.\ best particle vs.\ cluster mean]}.

Algorithm~\ref{alg:mcl} itself is generic and does not depend on the robot or its sensors.
What makes it specific to the Cyclops rover are the two models it calls:
\textsc{sample\_motion\_model} and \textsc{measurement\_model}. The next two sections
describe how each was built for the rover.

\section{Cyclops Rover Motion Model}
\label{sec:motion-model}
What specifically did we do 
\section{Cyclops Rover Measurement Model}
\label{sec:measurement-model}
Latency measured from laptop to be around 300ms

\section{Pathplanning}
\label{sec:navigation}

\red{some screenshots of different screen, map screen, real life, camera view}

%%% Local Variables: ***
%%% mode: latex ***
%%% TeX-master: "thesis.tex" ***
%%% End: ***
